%\documentclass[dvipdfmx,a4paper,12pt]{amsart}
\documentclass[dvipdfmx,a4paper,11pt]{report} %titleとe-mailをコメントアウトする．


%%% Packages %%%
\setlength{\lineskip}{0pt}

% --- 和文フォント設定 (ゴシックを使うなら) ---amsartを使う時はコメントアウト
\renewcommand{\kanjifamilydefault}{\gtdefault}
\usepackage{otf}          % min10を避けるため
\usepackage{pxrubrica}    % 和文ルビ

% --- 基本パッケージ ---
\usepackage{graphicx}
\usepackage[all]{xy}
\usepackage{wrapfig}
\usepackage{pgfplots}
\usepackage{color}
\usepackage[dvipsnames]{xcolor}

% --- 数学関連 ---
\usepackage{amsmath,amssymb,amsthm,amsfonts,mathtools}
\usepackage{amscd,dsfont,bigdelim,braket,physics,mathrsfs,bm}

% --- 書式・リスト関連 ---
\usepackage{latexsym}
\usepackage{setspace}
\usepackage{multirow}
\usepackage{enumerate}
\usepackage{enumitem}

% --- コメント・取消線など ---
\usepackage{comment}
\usepackage[normalem]{ulem} % \emph の下線化を抑止（cancelと共存）

% --- URL・文字コード ---
\usepackage{url}
% \usepackage[utf8]{inputenc}  % ← 使用しているエンジンがuplatexなら不要、pdflatexなら有効に

% --- showkeys（常に表示） ---
\usepackage{showkeys}
\renewcommand*{\showkeyslabelformat}[1]{%
 \fbox{\parbox{1.6cm}{\normalfont\tiny\sffamily#1\vspace{6mm}}}%
}
% --- hyperrefは最後に読み込む ---
\usepackage[dvipdfmx,colorlinks,linkcolor=blue,anchorcolor=blue,citecolor=blue]{hyperref}

%%% レイアウト調整 %%%
%%% レイアウト調整（geometryに統一） %%%
\usepackage[
  top=30mm,        % 上余白
  bottom=30mm,     % 下余白
  left=25mm,       % 左余白
  right=25mm,      % 右余白
  headheight=12pt, % ヘッダー高さ（必要なら）
  headsep=10mm,    % ヘッダーと本文の間
  footskip=32pt,   % 本文とフッターの距離
  includehead,     % ヘッダー分も高さに含める
  includefoot      % フッター分も高さに含める
]{geometry}

%%% 行間調整（適宜 1.2 などに変更） %%%
\usepackage{setspace}
\setstretch{1.1}

% --- 段落設定 --- 文字の行間ならここを変更する
\setlength{\parskip}{5pt}   % 段落間のスペース
\setlength{\parindent}{0pt}   % 段落先頭の字下げをなくす


%%% 追加（重複なし）パッケージ・設定 %%%

% --- 目次の体裁調整 ---
%\usepackage{tocloft}
%\renewcommand{\contentsname}{目次} % 日本語化
%\setlength{\cftbeforesecskip}{0pt}
%\setlength{\cftbeforesubsecskip}{0pt}
%\setlength{\cftbeforesubsubsecskip}{0pt}

% --- セクション見出しの体裁調整 ---
%\usepackage{titlesec}
%\titleformat*{\section}{\Large\bfseries}
%\titleformat*{\subsection}{\large\bfseries}
%\titlespacing*{\section}{0pt}{1.5ex plus .2ex minus .2ex}{0.8ex plus .1ex}
%\titlespacing*{\subsection}{0pt}{1.0ex plus .2ex minus .2ex}{0.5ex plus .1ex}

% --- ヘッダー／フッター設定 ---
%\usepackage{fancyhdr}
%\pagestyle{fancy}
%\fancyhf{}
%\rhead{岩井 雅崇}
%\lhead{大阪大学 数学専攻}
%\cfoot{\thepage}


% -- enumerate, itemize行間設定
\usepackage{enumitem} % デフォルト設定
\setlist[itemize]{itemsep=3pt, parsep=0pt}
\setlist[enumerate]{itemsep=3pt, parsep=0pt}
% "変更する際は右を使う" \setlength{\parskip}{0cm} % 段落間\setlength{\itemsep}{5pt} % 項目間

% --- tcolorbox 設定 ---%\begin{tcolorbox}[mybox]と使う
\usepackage[most]{tcolorbox}
\tcbuselibrary{breakable, skins, theorems}
\tcbset{
  mybox/.style={
    colback = white,
    colframe = green!35!black,
    fonttitle = \bfseries,
    breakable = true
  }
}

% --- TikZ 設定 ---
\usepackage{tikz}
\usetikzlibrary{positioning, arrows.meta, fit, calc, backgrounds}
\pgfdeclarelayer{background}
\pgfdeclarelayer{foreground}
\pgfsetlayers{background,main,foreground}

% --- footnote がページをまたがない設定 ---
\interfootnotelinepenalty=10000

% --- 目次に表示する階層の深さ ---
\setcounter{tocdepth}{2}

% --- 日本語目次---
\usepackage{pxjahyper}

%--newtheorem%--newcommand----

\newtheorem{thm}{Theorem}[section] 
\newtheorem{theo}[thm]{Theorem}
\newtheorem{cor}[thm]{Corollary}
\newtheorem{prop}[thm]{Proposition}
\newtheorem{conj}[thm]{Conjecture}
\newtheorem*{mainthm}{Theorem}
\newtheorem{deflem}[thm]{Definition-Lemma}
\newtheorem{lem}[thm]{Lemma}
\theoremstyle{definition} 
\newtheorem{defn}[thm]{Definition}
\newtheorem{propdefn}[thm]{Proposition-Definition} 
\newtheorem{lemdefn}[thm]{Lemma-Definition} 
\newtheorem{thmdefn}[thm]{Theorem-Definition} 
\newtheorem{eg}[thm]{Example} 
\newtheorem{ex}[thm]{Example} 
\newtheorem{ques}[thm]{Question}
\newtheorem{remin}[thm]{Reminder}
\theoremstyle{remark}
\newtheorem{rem}[thm]{Remark}
\newtheorem{setup}[thm]{Setup}
\newtheorem{obs}[thm]{Observation}
\newtheorem{notation}[thm]{Notation}
\newtheorem{cl}{Claim}
\newtheorem{claim}[thm]{Claim}
\newtheorem{assup}[thm]{Assumption}
\newtheorem{step}{Step}
\newtheorem*{clproof}{Proof of Claim}
\newtheorem{cln}[thm]{Claim}
\newtheorem*{ack}{Acknowledgements} 
\numberwithin{equation}{section}
\newtheorem{case}{Case}

\newcommand{\rk}[0]{\operatorname{rk}}
\newcommand{\supp}[0]{\operatorname{Supp}}
\newcommand{\Rad}[0]{\operatorname{Rad}}
\newcommand{\Sha}[0]{\operatorname{Sha}}
\newcommand{\sha}[0]{\operatorname{sha}}
\newcommand{\eend}[0]{\operatorname{End}}
\newcommand{\codim}[0]{\operatorname{codim}}
\newcommand{\nd}[0]{\operatorname{nd}}
\renewcommand{\rank}[0]{\operatorname{rank}}
\newcommand{\degree}[0]{\operatorname{deg}}
\newcommand{\Exc}[0]{\operatorname{Exc}}
\newcommand{\pr}{{\rm pr}}
\newcommand{\id}{{\rm id}}
\newcommand{\Sym}{{\rm Sym}}
\newcommand{\End}[0]{\operatorname{End}}
\newcommand{\Coker}[0]{\operatorname{Coker}}

\newcommand{\Supp}{{\rm Supp}}
\newcommand{\Hom}[0]{\mathscr{H}\!\textit{om}}
\newcommand{\Ext}[0]{\mathscr{E}\!\textit{xt}}
\newcommand{\GL}[0]{\operatorname{GL}}
\newcommand{\SheafHom[1]}{\mathscr{H}\!\!\!\text{\calligra om}_{\,{#1}}}
\newcommand{\PGL}[0]{\mathbb{P}\GL(r,\C)}

\newcommand{\Alb}{{\rm Alb}}
\newcommand{\verti}{{\rm vert}}
\newcommand{\hor}{{\rm hor}}
\newcommand{\univ}{{\rm univ}}
\newcommand{\Tor}{{\rm tor}}
\newcommand{\shaf}{\mathrm{sha}}
\newcommand{\Shaf}{\mathrm{Sha}}
\newcommand{\reg}{{\rm{reg}}}
\newcommand{\sing}{{\rm{sing}}}
\newcommand{\qt}{{\rm{qt}}}
\newcommand\sO{{\mathcal O}}
\newcommand{\Div}[0]{\operatorname{div}}
\newcommand{\ddbar}{dd^c}
\newcommand{\cV}{\mathcal{V}}
\newcommand{\deldel}{\sqrt{-1}\partial \overline{\partial}}
\newcommand{\dbar}{\overline{\partial}}
\newcommand{\I}[1]{\mathcal{I}(#1)}
\newcommand{\Aut}[1]{\mathrm{Aut}(#1)}
\newcommand{\Ker}[1]{\mathrm{Ker}(#1)}
\newcommand{\Image}[1]{\mathrm{Im}(#1)}
\DeclareMathOperator{\Ric}{Ric}
\DeclareMathOperator{\Vol}{Vol}
 \newcommand{\pdrv}[2]{\frac{\partial #1}{\partial #2}}
 \newcommand{\drv}[2]{\frac{d #1}{d#2}}
  \newcommand{\ppdrv}[3]{\frac{\partial #1}{\partial #2 \partial #3}}
\newcommand{\underalign}[2]{\quad \underset{\mathclap{\strut #1}}{#2}\quad}
\newcommand{\polar}{\beta}
  
\newcommand{\R}{\mathbb{R}}
\newcommand{\Z}{\mathbb{Z}}
\newcommand{\N}{\mathbb{Z}_+}
\newcommand{\C}{\mathbb{C}}
\newcommand{\Q}{\mathbb{Q}}
\newcommand{\D}{\mathbb{D}}
\newcommand{\mP}{\mathbb{P}}
\newcommand{\mO}{\mathcal{O}}
\newcommand{\B}{\mathds{B}}
\newcommand{\tl}{\hspace{-0.8ex}<\hspace{-0.8ex}}
\renewcommand{\tr}{\hspace{-0.8ex}>}

\newcommand{\xb}[1]{\textcolor{blue}{#1}}
\newcommand{\xr}[1]{\textcolor{red}{#1}}
\newcommand{\xm}[1]{\textcolor{magenta}{#1}}

\newcommand{\illegible}[1]{\textcolor{red}{[ILLEGIBLE: #1]}}

\title{カレントセミナーのノート} 
\author{岩井雅崇}
%\address{Department of Mathematics, Graduate School of Science, Osaka City University 3-3-138, Sugimoto, Sumiyoshi-ku Osaka, 558-8585Japan} 
%\email{{\tt masataka.math@gmail.com}}
%\email{{\tt masataka.math@gmail.com, masataka@sci.osaka-cu.ac.jp}}
\date{\today, version 0.01}

\renewcommand{\thefootnote}{\arabic{footnote}}

\baselineskip = 15pt
\footskip = 32pt

\begin{document}
\begin{thebibliography}{n}
\bibitem[Rud]{Rud}
W. Rudin. \textit{Functional analysis.} 2nd edn. International Series in Pure and Applied Mathematics. McGraw-Hill, Inc., New York. (1991.)
\bibitem[Rud 2]{Rud2}
W. Rudin. \textit{Real and complex analysis. } 3rd ed. New York, NY: McGraw-Hill 
\bibitem[NO]{NO}
J. Noguchi, T.Ochiai \textit{Geometric Function Theory in Several Complex Variables} Translations of Mathematical Monographs
Volume: 80; 1990; 282 pp
\end{thebibliography}
 
\cite[Chapter 3]{NO}を参考にしている. 

\section{}

\begin{tcolorbox}[mybox]
\begin{prop}
 \(M\) 向きづけ可能 \(C^\infty\)多様体. 
  \(\dim_{\mathbb R} M = m\),  \(0 \le p \le m\)とする. 
\[
[\,\cdot\,] : \mathcal{L}^{m-p}_{\mathrm{loc}}(M) \to \mathcal{K}'_p(M)\ (\subset \mathcal{D}'_p(M))
\quad
\omega \longmapsto \left(\varphi \mapsto \int_M \omega \wedge \varphi\right)
\]
は全射となる. 
\end{prop}
\end{tcolorbox}

\begin{rem}
 \(\exists f:\mathbb{R}\to\mathbb{R}\) injective \(C^\infty\)関数と \(\exists A \subset \mathbb{R}\) Lebesgue可測なもので. 
\(f^{-1}(A)\) がLebesgue可測にならないものがある.

\(\exists f:[0,1]\to[0,1]\) 同相写像, \(\exists A \subset \mathbb{R}\) Lebesgue可測なもので. 
\(f^{-1}(A)\) がLebesgue可測にならないものがある.
\end{rem}
\begin{ques}
正則写像やalgebraic mapの場合はどうなるか？
\end{ques}

\begin{proof}
% [page 1]
 \([\,\cdot\,]: \mathcal{L}^{m-p}_{\mathrm{loc}}(\bullet) \to \mathcal{K}'_p(\bullet)\)がsheafの射なので,  \(M = U \subset \mathbb{R}^m\)という開集合にして良い. 
さらに \(\omega = \sum_J \omega_J dx^J \in L^{m-k}_{\mathrm{loc}}(U)\)とする．
\footnote{ただし$ L^{m-k}$はalmost everywhere で同一視する前のもので, 係数は関数とする. (つまり同一視する前のもの)}

任意の$\varphi \in \mathcal{K}^k(U)$について, \(\int_U \omega \wedge \varphi = 0\)であると仮定する. 
示すことは任意の添字$J$について,  \(\omega_J = 0\) である. 

任意の \(J=(i_1,\cdots,i_{m-k})\)について, 
 \(J^\circ := (j_1,\cdots,j_k)\) を
\[
\{i_1,\cdots,i_{m-k},\, j_1,\cdots,j_k\}=\{1,\cdots,m\}.
\]
となるようにとる. 任意の\(\varphi \in \mathcal{K}(U)\)について, 
\[
0=
\int_U \omega \wedge \underbrace{\varphi\, dx^{J^\circ}}_{\in \mathcal{K}^k(U)}
   = \int_U \omega_J \varphi\, dx^J \wedge dx^{J^\circ}.
\]
よって, 任意の\(\forall \varphi \in \mathcal{K}(U)\)について, 
\(\int_U \omega_J \varphi\, dm = 0\)である. ここで$dm$をLesbegue 測度とする. 
よって次の補題からいえた. 
\end{proof}

\begin{tcolorbox}[mybox]
\begin{lem}
\(U \subset \mathbb{R}^m\) 開集合, \(f \in L_{\mathrm{loc}}(U)\)とする. 
今任意の$\varphi \in \mathcal{K}(U)$について. 
\(
\int_U f\varphi\, dm = 0 
\)
が成り立つならば,  測度\(m\)についてalmost everywhereで\(f=0\)である.
\end{lem}
\end{tcolorbox}
$ \mathcal{K}(U)$
は$U$上のコンパクトサポートを持つ連続$\C$値関数の集合である.
\begin{proof}

[Step 1.] 
\(f \in L^1(U)\) (つまり\(\int_U |f|\,dm < \infty\))であることを仮定して良い. 

なぜならば \(U=\bigcup_{i\geq 1} U_i\) 可算個の開被覆で
\(\overline{U_i}\) がコンパクトとなるものをとる. 
 \(\{\varphi_i\}_{i\geq 1}\) を１の分割とする. (Supp \(\varphi_i \subset U_i\)となる)

すると\(\varphi_i f \in L^1(U)\) かつ \(\forall \varphi \in \mathcal{K}(U)\)について
\[
\int_U (\varphi_i f)\varphi\,dm
=
\int_U f \underbrace{(\varphi_i\cdot \varphi)}_{\in \mathcal{K}(U)}\,dm
=
0
\]
よって, almost everywhereで$\varphi_i f=0$となるので, 
\[
f=\sum_{i\geq 1}\varphi_i f = 0 \ \text{a.e.}
\]

[Step 2] 
\(\quad f \in L^1(U)\) の場合に示す. 

%Idea.\[\mathcal{K}(U) \to \mathbb{C};\ \varphi \mapsto \int_U f\varphi\,dm\]
%\textcolor{red}{[囲み内注記]}
%これは supp narrow にいて bdd \textcolor{red}{\illegible{英字略記}}. Riesz rep
%\textcolor{red}{\illegible{日本語文}}

任意のLesbegue可測集合 \(E \subset U\) について, 
\[
\lambda(E)=\int_E f\,dm.
\]
とする. すると\(\lambda\)はcomplex measureになる.\footnote{$\lambda$がcomplex measureとは
$\lambda:\ \{\text{Leb meas set}\}\to \mathbb{C}$であって, 完全加法性
$\lambda\Bigl(\bigsqcup_{i=1}^{\infty} A_i\Bigr)
=
\sum_{i=1}^{\infty}\lambda(A_i)$が成り立つこと. ただし右辺の収束は絶対収束を要求する}

さて$\lambda$の total variationを
Lesbegue可測集合\(E\subset U\)について
\[
|\lambda|(E):=\sup\left\{\sum_{i=1}^{\infty}|\lambda(E_i)|\ \middle|\ E=\bigsqcup_{i=1}^{\infty}E_i,\ E_i\subset U \text{Lesbegue可測}\right\}
\]
とする. 
\cite[Chapter 6]{Rud2}より,  \(|\lambda|\) は finite positive measure, つまり$|\lambda|(U)<\infty$となる測度である. (この測度$|\lambda|$はtotal variation measureと呼ばれる. )

\cite[Theorem 6.12]{Rud2}よりRadon-Nikodymの定理から
ある\( h \in L^1(U,|\lambda|)\)であって, 
任意の$x \in U$について\(
|h(x)|=1\)かつ
任意のLesbegue可測集合$E\subset U$について
\[
\lambda(E)=\int_E h\,d|\lambda|.
\]
となるものが存在する. \footnote{このことを$d\lambda = h\,d|\lambda|$
と表し. \(\lambda\)のpolar expressionともいう}

このとき任意の\(\varphi \in \mathcal{K}(U)\)について. 
 \[\displaystyle \underbrace{\int_U \varphi h\,d|\lambda|}_{=:\int_U \varphi\,d\lambda} = \int_U \varphi f\,dm \underset{\text{仮定}}{=} 0\]
これを示す.

まず$E$をLesbegue可測集合で$\varphi=\chi_E$, つまり
$$\varphi=\chi_E,\qquad
\chi_E(x)=
\begin{cases}
1^\circ & x\in E\\
0 & x\notin E
\end{cases}
$$
である場合は
\[
\int_U \chi_E h\,d|\lambda|
\underset{\text{$\chi_E$の定義}}{=}
\int_E h\,d|\lambda|
\underset{\text{定義}}{=}
\lambda(E)
\underset{\text{定義}}{=}
\int_E f\,dm
\underset{\text{$\chi_E$の定義}}{=}
\int_U \chi_E f\,dm.
\]
であるのでOK.
\( \varphi\)が単関数(support関数$\chi_E$の有限な線型結合となる関数)の時もOK.

$ \varphi$が0以上の有界関数の時, 
単関数近似
\(s_i \to \varphi \)をとる. (ただし$s_i \le s_{i+1} \le \varphi $となるようにとる)
すると
\[
\int_U s_i h\,d|\lambda| \xrightarrow[i\to\infty]{} \int_U \varphi h\,d|\lambda|
\quad 
\int_U s_i f\,dm \xrightarrow[i\to\infty]{} \int_U \varphi f\,dm
\]
である, そして
\[
|s_i h|\leq \underbrace{|\varphi|}_{\text{bounded}}|h|,\qquad h\in L^1 \text{ w.r.t. } |\lambda|
\quad \text{and} \quad
|s_i f|\leq |\varphi||f| \in L^1 \text{ w.r.t. } m
\]
であるので, Lesbegue収束定理より$\int_U \varphi h\,d|\lambda|= \int_U \varphi f\,dm$となる. 
以上より$\varphi \in \mathcal{K}(U)$の時でもOKである. 

さて\(|\lambda|(U)=0\)ならば \(|\lambda|=0\) であることを示す. 
もしこれがいえたら\(\lambda=0\)が言えて,  任意のLesbegue可測集合 \( E\subset U\) について
\[
0\underset{\lambda=0}{=}\lambda(E)
\underset{\text{定義}}{=}\int_E f\,dm.
\]
となり, $f=0$ a.e.が言える.
今
\begin{align*}
|\lambda|(U)
&=\int_U d|\lambda|
\underset{|h|=1}{=}
\int_U \overline{h}\,h\,d|\lambda| \\
&\underset{\forall \varphi\in \mathcal{K}(U)}{=}\int_U (\overline{h}-\varphi)h\,d|\lambda|
\underset{}{=}\int_U \operatorname{Re}\bigl((\overline{h}-\varphi)h\bigr)\,d|\lambda| \\
&\underset{}{\le}
\int_U |(\overline{h}-\varphi)h|\,d|\lambda|\\
&\underset{|h|=1}{=}\int_U |\overline{h}-\varphi|\,d|\lambda|
\end{align*}
ここで$\int_U |\overline{h}-\varphi|\,d|\lambda|$は
\(\overline{h}\) と \(\varphi\) の \(L^1(U,|\lambda|)\)-距離である.

よって下のclaimより, $\overline{h} \in L^1(U,|\lambda|)$について
$\varphi_i \to \overline{h}$となる$\varphi_i \in \mathcal{K}(U)$をとれば, 
$|\lambda|(U)=0$がいえる. 

\begin{claim}
\(\mathcal{K}(U)\subset L^1(U,|\lambda|)\) は
 \(L^1\)-normでdenseである. 
\end{claim}
次の定理を思い出す. 
\begin{tcolorbox}[mybox]
\begin{thm}\cite[Theorem 3.14]{Rud2}
\(X\) locally compact Hausdorff,  \(\mathcal{M}\) $X$の\(\sigma\)-algebra, 
\(\mu\)を\(\mathcal{M}\)上のpositive measureとする. 

 \((X,\mathcal{M},\mu)\) が次を満たすとする.
\begin{itemize}
\item[(0)] \(\mathcal{M}\) は Borel algebraを含む. (つまり$\mathcal{M}$が$X$の全ての開集合を含む)
\item[(1)] 任意のコンパクト集合 \( K\subset X\) について \(\mu(K)<\infty\)
\item[(2)] (外部正則性) 任意の\( E\in\mathcal{M}\)について$\mu(E)=\inf\{\mu(V)\mid E\subset V\underset{\text{open}}{\subset} X \}\)
\item[(3)] (内部正則性) 任意の \(E\in\mathcal{M}\), \(E\) が開集合または\(\mu(E)<\infty\)ならば
\(
\mu(E)=\sup\{\mu(K)\mid K\underset{\text{cpt}}{\subset} E\}
\)
\item[(4)](完備性) 任意の\(E\in\mathcal{M},\ A\subset E,\ \mu(E)=0\)について, $A\in\mathcal{M}$
\end{itemize}
% [page 1]

この時 \(\{f:X\to\mathbb{C}\ \text{conti} \mid \operatorname{supp} f \text{ cpt}\}\subset L^1(X,\mu)\)
は\(L^1\)-normで
denseである. 
実は$L^p,\ 1\leq p<\infty$でも成り立つ. 
\end{thm}
\end{tcolorbox}

\(|\lambda|\) が定理の仮定を満たすことをみる. 
f\(|\lambda|\)は有限測度で, 
$ (0)\sim (3)$は OK(多少議論は必要, 有限Borel測度の正則性)
(4) については \((U,\text{Lesbegue可測集合},|\lambda|)\) の
測度空間としての完備かをとれば, 
\((0)\sim(3)\) を満たしながら, (4)も成立する.
この時$L^1$空間が大きくなるだけである. よってclaimが成り立つ. 
\end{proof}


\end{document}